\chapter*{Abstract}
\addcontentsline{toc}{chapter}{Abstract}

This Integrative Project describes the design, construction, and commissioning of a smart vertical urban farm prototype. The system's objective is to automate indoor farming by autonomously monitoring and controlling the main environmental variables: temperature, air humidity, light intensity, and soil moisture. For its development, a multi-level physical structure adapted for traditional soil-based cultivation was built. The system features a central processing unit responsible for receiving sensor data and activating the various control devices. These include artificial lighting panels, a localized irrigation network with water recovery, and an automated ventilation system. The overall assembly was completed with custom-made mechanical parts to ensure the proper support of all components. On the software side, a control logic was programmed to maintain climate and irrigation within the necessary parameters for the plants at all times. Furthermore, the equipment features an internet connection to send data to the cloud, making it possible to monitor the crop's status remotely and in real time, thus applying the Internet of Things (IoT) concept.

\vspace{1cm}

\noindent\textbf{Thematic area:} Digital Electronics and Control Systems

\vspace{0.5cm}

\noindent\textbf{Subjects:} Digital Electronics 3, Industrial Electronics, Analog Electronics 1, and Control Systems 1

\vspace{0.5cm}

\noindent\textbf{Key words:} farm, cultivation, substrate, farms, smart cities, vertical